\section*{Bab \ref{ch:b1:continuity} --- Limit dan Kekontinuan}

\begin{solution}{exo:b1:continuity:1}
Ambillah $u_n = \frac{1}{2\pi n + \pi/2}$ dan $v_n = \frac{1}{2\pi n}$:
keduanya menuju $0^+$, namun $\sin\frac{1}{u_n} = 1$ dan
$\sin\frac{1}{v_n} = 0$. Jadi dua barisan, dengan dua limit peta yang
berbeda: sehingga menurut \cref{thm:b1:continuity:seqchar}, tak ada limit di $0^+$.

Adapun $x \sin\frac1x$: terapit oleh $\abs{x\sin\frac1x} \leq \abs x \to 0$,
sehingga limitnya di $0$ ada dan sama dengan $0$.
\end{solution}

\begin{solution}{exo:b1:continuity:2}
Fungsi $f = \lfloor\cdot\rfloor$ \omterm{def:b1:continuity:continuous}{kontinu} pada $\R \setminus \Z$
(karena konstan secara lokal) dan takkontinu di setiap $n \in \Z$: dengan limit kirinya
$n - 1$, sedangkan nilainya $n$.

$g(x) = x - \lfloor x\rfloor$: dengan titik ketakkontinuan yang sama (karena
identitasnya \omterm{def:b1:continuity:continuous}{kontinu}, sehingga $g$ mewarisi lompatan $f$); di $n \in
\Z$, limit kirinya $1 \neq 0 = g(n)$.

$h(x) = \lfloor x\rfloor + (x - \lfloor x\rfloor)^2$: pada $\intco{n}{n
+ 1}$, $h(x) = n + (x - n)^2$, yang \omterm{def:b1:continuity:continuous}{kontinu} di sana; sedangkan di $x = n$, limit
kirinya $(n-1) + 1 = n = h(n)$: jadi lompatannya saling meniadakan. Sehingga $h$ \omterm{def:b1:continuity:continuous}{kontinu}
pada $\R$ (dan naik tegas).
\end{solution}

\begin{solution}{exo:b1:continuity:3}
$P(x) = x^5 - 3x + 1$: di sini $P(-2) = -32 + 6 + 1 = -25 < 0$; $P(0) = 1 >
0$; $P(1) = -1 < 0$; $P(2) = 27 > 0$. Jadi ada tiga perubahan tanda pada
ruas yang lepas $\intcc{-2}{0}$, $\intcc{0}{1}$, $\intcc{1}{2}$:
sehingga menurut \cref{thm:b1:continuity:ivt}, sekurang-kurangnya ada tiga akar. (Karena berderajat
$5$, $P$ mempunyai paling banyak lima; dan telaah variasinya akan menunjukkan
tepat tiga.)
\end{solution}

\begin{solution}{exo:b1:continuity:4}
Misalkan $P = a_{2m+1}X^{2m+1} + \dots$ dengan $a_{2m+1} > 0$ (kalau tidak gantilah
$P$ dengan $-P$). Dengan memfaktorkan suku dominannya, $P(x) = a_{2m+1}
x^{2m+1}\bigl(1 + o(1)\bigr)$ ketika $x \to \pm\infty$: sehingga $P(x) \to
+\infty$ di $+\infty$ dan $-\infty$ di $-\infty$. Pilihlah $a$ dengan
$P(a) < 0$ dan $b$ dengan $P(b) > 0$: maka teorema nilai antara pada
$\intcc{a}{b}$ menyediakan sebuah akar.
\end{solution}

\begin{solution}{exo:b1:continuity:5}
Misalkan $g(x) = f(x) - x$, yang \omterm{def:b1:continuity:continuous}{kontinu} pada $\intcc{0}{1}$. Karena $f$ memetakan
ke dalam $\intcc{0}{1}$: maka $g(0) = f(0) \geq 0$ dan $g(1) = f(1) - 1 \leq
0$. Menurut \cref{thm:b1:continuity:ivt}, $g(c) = 0$ untuk suatu $c$: yaitu
sebuah titik tetap.

Keperluan hipotesisnya: pada $\intcc{0}{1}$, \omterm{def:b1:logic:map}{pemetaan} takkontinu
$f(x) = 1$ untuk $x \leq \frac12$, $f(x) = 0$ untuk $x > \frac12$ tak mempunyai
titik tetap; sedangkan pada \omterm{prop:b1:reals:intervals}{selang} $\intoo{0}{1}$ (yang bukan ruas), $f(x) =
\frac x2$ \omterm{def:b1:continuity:continuous}{kontinu} ke dalam $\intoo{0}{1}$ tanpa titik tetap
(karena calonnya $0$ tak ada); dan pada $\R$, $f(x) = x + 1$.
\end{solution}

\begin{solution}{exo:b1:continuity:6}
Fungsi $g(x) = \cos x - x$ \omterm{def:b1:continuity:continuous}{kontinu}, dengan $g(0) = 1 > 0$, $g(1) = \cos 1 -
1 < 0$: sehingga sebuah penyelesaian ada di $\intoo{0}{1}$
(\cref{thm:b1:continuity:ivt}). Ketunggalannya: karena $g$ turun
tegas pada $\R$ --- untuk $x \leq 0$, $g(x) \geq 1 - x > 0$ toh tak mempunyai
nol; dan $g'(x) = -\sin x - 1 \leq 0$ dengan kesamaannya hanya di
titik yang terasing ($x \equiv -\frac\pi2 \bmod 2\pi$), sehingga $g$ turun
tegas (\cref{ch:b1:derivative}; atau sebagai alternatifnya: pada
$\intcc{0}{1}$, $\cos$ turun tegas dan $-x$ pun demikian, sehingga $g$
juga). Adapun fungsi yang monoton tegas lenyap paling banyak sekali.
\end{solution}

\begin{solution}{exo:b1:continuity:7}
Tetapkan $M = f(0) + 1$. Ada $A > 0$ dengan $f(x) \geq M$ untuk $\abs x
\geq A$ (menurut definisi kedua limit tak hingganya; ambillah ambang yang
lebih besar). Pada ruas $\intcc{-A}{A}$, teorema nilai ekstrem
(\cref{thm:b1:continuity:evt}) menyediakan $c$ dengan $f(c) =
\inf_{\intcc{-A}{A}} f \leq f(0)$. Untuk $\abs x \geq A$: $f(x) \geq M
> f(0) \geq f(c)$. Jadi $f(c)$ merupakan minimum semestanya.
\end{solution}

\begin{solution}{exo:b1:continuity:8}
Pada ruas $\intcc{0}{T}$, $f$ terbatas dan mencapai batasnya
(\cref{thm:b1:continuity:evt}); lalu menurut keperiodikannya, itulah batasnya
pada seluruh $\R$, yang tetap tercapai.

Misalkan $g(x) = f(x + \frac T2) - f(x)$, yang \omterm{def:b1:continuity:continuous}{kontinu}. Maka
\[
g(0) + g\bigl(\tfrac T2\bigr)
= \bigl(f(\tfrac T2) - f(0)\bigr) + \bigl(f(T) - f(\tfrac T2)\bigr)
= f(T) - f(0) = 0 :
\]
jadi $g(0)$ dan $g(\frac T2)$ bertanda berlawanan (atau salah satunya lenyap), sehingga
teorema nilai antara pada $\intcc{0}{T/2}$ memberikan $c$ dengan
$g(c) = 0$, yakni $f(c + \frac T2) = f(c)$.
\end{solution}

\begin{solution}{exo:b1:continuity:9}
Pertama ketaksamaannya: untuk $0 \leq y \leq x$,
\[
\bigl(\sqrt y + \sqrt{x - y}\bigr)^2 = x + 2\sqrt{y(x-y)} \geq x,
\]
sehingga $\sqrt x \leq \sqrt y + \sqrt{x - y}$, yakni $\sqrt x - \sqrt y
\leq \sqrt{x - y}$. Jadi $\abs{\sqrt x - \sqrt y} \leq
\sqrt{\abs{x - y}}$ untuk setiap $x, y \geq 0$.

\omterm{def:b1:continuity:uniform}{Kekontinuan seragamnya}: diberikan $\varepsilon > 0$, ambillah $\delta =
\varepsilon^2$; maka $\abs{x - y} \leq \delta$ mengakibatkan $\abs{\sqrt x
- \sqrt y} \leq \sqrt\delta = \varepsilon$.

Tak Lipschitz di dekat $0$: karena $\frac{\sqrt x - \sqrt 0}{x - 0} =
\frac{1}{\sqrt x} \to +\infty$ ketika $x \to 0^+$, sehingga tak ada konstanta $k$
yang dapat mendominasi semua hasil bagi selisihnya.
\end{solution}

\begin{solution}{exo:b1:continuity:10}
Andaikan $f$ \omterm{def:b1:logic:inj}{injektif}, \omterm{def:b1:continuity:continuous}{kontinu}, dan tak monoton tegas.
Maka ada $a < b < c$ di $I$ dengan $f(b)$ tak berada di antara $f(a)$ dan
$f(c)$ --- memang, seandainya untuk \emph{setiap} tripelnya nilai tengahnya berada
di antara nilai luarnya, maka $f$ akan monoton (bandingkanlah sebarang dua
pasangnya; lewat pemeriksaan kasus yang singkat). Katakanlah $f(b) > \max(f(a), f(c))$ (adapun kasus
yang lain simetris, gantilah $f$ dengan $-f$). Pilihlah $v$ dengan $\max(f(a),
f(c)) < v < f(b)$. Menurut teorema nilai antara yang diterapkan pada
$\intcc{a}{b}$ dan pada $\intcc{b}{c}$, ada $u_1 \in
\intoo{a}{b}$ dan $u_2 \in \intoo{b}{c}$ dengan $f(u_1) = v = f(u_2)$:
yaitu dua titik berbeda dengan peta yang sama, yang bertentangan dengan keinjektifannya.
\end{solution}

\begin{solution}{exo:b1:continuity:11}
Dari $f(0) = f(0+0) = 2f(0)$ diperoleh $f(0) = 0$; dan $f(-x) = -f(x)$ dari $0 =
f(x - x)$. Tetapkan $\alpha = f(1)$. Lewat induksi: $f(n) = n\alpha$ untuk $n
\in \N$, lalu untuk $n \in \Z$ menurut keganjilannya. Untuk $q \in \N^*$:
$q\,f(\frac pq) = f(p) = p\alpha$ (dengan menjumlahkan $\frac pq$ pada dirinya $q$
kali), sehingga $f(\frac pq) = \alpha\frac pq$: jadi $f = \alpha\,
\mathrm{id}$ pada $\Q$.

Sekarang misalkan $x \in \R$ dan $(r_n)$ barisan bilangan rasional dengan $r_n \to
x$ (menurut \omterm{def:b1:topology:dense}{kepadatannya}, \cref{thm:b1:reals:density}, yang diterapkan pada \omterm{prop:b1:reals:intervals}{selang}
yang bersarang; atau $r_n = \frac{\lfloor nx\rfloor}{n}$). Lewat \omterm{def:b1:continuity:continuous}{kekontinuannya}:
$f(x) = \lim f(r_n) = \lim \alpha r_n = \alpha x$.
\end{solution}

\begin{solution}{exo:b1:continuity:12}
Misalkan $\varepsilon > 0$. Menurut limitnya di $+\infty$, ada $A$ dengan
$\abs{f(x) - \ell} \leq \frac\varepsilon2$ untuk $x \geq A$; sehingga untuk
$x, y \geq A$: $\abs{f(x) - f(y)} \leq \varepsilon$ (tanpa memerlukan
kedekatan).

Pada ruas $\intcc{0}{A + 1}$, teorema Heine
(\cref{thm:b1:continuity:heine}) memberikan $\delta_0 > 0$ bagi
$\varepsilon$ ini; lalu tetapkan $\delta = \min(\delta_0, 1)$.

Sekarang ambillah sebarang $x, y \geq 0$ dengan $\abs{x - y} \leq \delta$, katakanlah $x
\leq y$. Jika $y \leq A + 1$: maka keduanya terletak di ruasnya, dan $\delta_0$
dari Heine berlaku. Kalau tidak, $y > A + 1$, lalu $x \geq y - 1 >
A$: sehingga keduanya terletak di $\intco{A}{+\infty}$, yang di situ argumen limitnya
berlaku. Pada kedua kasusnya $\abs{f(x) - f(y)} \leq \varepsilon$: jadi \omterm{def:b1:continuity:uniform}{kontinu
seragam}.
\end{solution}

\begin{solution}{pb:b1:continuity:1}
\textbf{1.} Untuk $n \in \N$: $f(nx) = n f(x)$ lewat induksi
(karena $f((n+1)x) = f(nx) + f(x)$). Lalu $f(0) = 2f(0)$ memberikan $f(0) =
0$, dan $0 = f(x - x) = f(x) + f(-x)$ memberikan keganjilannya, sehingga $f(nx) =
nf(x)$ untuk $n \in \Z$. Untuk $r = \frac pq$: $q\,f\bigl(\tfrac
pq x\bigr) = f(px) = p\,f(x)$, sehingga $f(rx) = r f(x)$: jadi $f$ bersifat
linear atas $\Q$.

\textbf{2.} Keaditifannya memberikan, untuk setiap $x$ dan $h$:
\[
f(x + h) - f(x) = f(h) = f(x_0 + h) - f(x_0) .
\]
Ketika $h \to 0$, ruas kanannya menuju $0$ menurut \omterm{def:b1:continuity:continuous}{kekontinuannya} di
$x_0$; sehingga $f(x + h) \to f(x)$: jadi \omterm{def:b1:continuity:continuous}{kontinu} di setiap $x$. Lalu
\cref{exo:b1:continuity:11} menghasilkan $f(x) = f(1)\,x$.

\textbf{3.} Dua fungsi aditif yang \omterm{def:b1:continuity:continuous}{kontinu} berbentuk
$cx$ dan $c'x$ (pertanyaan 2); lalu jika keduanya bersesuaian pada \omterm{def:b1:structures:subgroup}{subgrup} \omterm{def:b1:topology:dense}{padat}
$\Z + \sqrt2\,\Z$ (\cref{exo:b1:reals:9}), maka $c\,g = c'g$ untuk
suatu $g \neq 0$ di dalamnya: sehingga $c = c'$, jadi fungsinya sama. Adapun tanpa
\omterm{def:b1:continuity:continuous}{kekontinuan}: kelinearan atas $\Q$ hanya mengikat nilainya pada
kombinasi $\Q$-nya, dan $1, \sqrt2$ bebas atas $\Q$ (karena $\sqrt2
\notin \Q$), sehingga $f(1) = 0$, $f(\sqrt2) = 1$ tetap taat asas dan
memaksa, pada \omterm{def:b1:structures:subgroup}{subgrupnya},
\[
f(m + n\sqrt2) = m\,f(1) + n\,f(\sqrt2) = n .
\]

\textbf{4.} Untuk $t \in \intcc{0}{\ell}$: $a + t \in
\intcc{a}{b}$, sehingga $f(t) = f(a + t) - f(a) \leq M - f(a) =: M'$.

\textbf{5.} Untuk $t \in \intcc{0}{\ell}$, juga $\ell - t \in
\intcc{0}{\ell}$, dan keaditifannya memberikan $f(t) = f(\ell) - f(\ell
- t) \geq f(\ell) - M'$. Jadi $\abs{f} \leq C$ pada
$\intcc{0}{\ell}$ dengan $C = \max\bigl(\abs{M'}, \abs{f(\ell) -
M'}\bigr)$.

\textbf{6.} Untuk $t \in \intcc{0}{\ell/n}$: $nt \in
\intcc{0}{\ell}$ dan $f(t) = \frac{f(nt)}{n}$ (pertanyaan 1), sehingga
$\abs{f(t)} \leq \frac Cn$. Diberikan $\varepsilon > 0$, pilihlah $n >
\frac C\varepsilon$: maka untuk $0 \leq h \leq \frac\ell n$,
$\abs{f(h)} \leq \varepsilon$, dan untuk $h$ yang negatif pakailah $f(h) =
-f(-h)$. Jadi $f(h) \to 0 = f(0)$ ketika $h \to 0$: yakni \omterm{def:b1:continuity:continuous}{kontinu} di $0$,
sehingga di mana-mana (pertanyaan 2), jadi $f(x) = f(1)x$.

\textbf{7.} Jika $f$ tidak turun pada $\intcc{a}{b}$, maka
$f(a) \leq f(x) \leq f(b)$ di sana: jadi terbatas, sehingga linear menurut pertanyaan
6, yakni $f(x) = cx$; lalu $c(b - a) = f(b) - f(a) \geq 0$ memaksa $c
\geq 0$.

\textbf{8.} (a)$\Rightarrow$(b)$\Rightarrow$(c): sepele.
(c)$\Rightarrow$(a): pertanyaan 2. (a)$\Rightarrow$(d): karena fungsi
linear bersifat monoton di mana-mana. (d)$\Rightarrow$(e): karena fungsi monoton
pada $\intcc{a}{b}$ terbatas di sana oleh nilai titik
ujungnya. (e)$\Rightarrow$(a): pertanyaan 4--6. Jadi kelima
\omterm{def:b1:logic:statement}{pernyataannya} setara --- itulah tangga keteraturannya.

\textbf{9.} Pertanyaan 4 hanya memakai batas atasnya $M$; sedangkan pertanyaan 5
\emph{menurunkan} batas bawahnya dari batas atasnya lewat
pencerminan $f(t) = f(\ell) - f(\ell - t)$; lalu pertanyaan 6 berjalan di atas
$\abs f \leq C$. Jadi: aditif dan terbatas \emph{di atas} pada satu
\omterm{prop:b1:reals:intervals}{selang} yang taktermerosot sudah mengakibatkan linear. Adapun untuk batas
bawahnya, terapkanlah ini pada $-f$ (yang aditif dan terbatas di atas). Sehingga anak tangga
(e) yang terkuat: bahwa satu batas sepihak pada satu \omterm{prop:b1:reals:intervals}{selang} sudah mencukupi.

\textbf{10.} Seandainya $\frac{f(x)}{x}$ merupakan satu konstanta $c$ untuk setiap
$x \neq 0$, maka $f$ akan linear; jadi ada $u, v$ yang taknol
dengan $\frac{f(u)}u \neq \frac{f(v)}v$, yakni $u f(v) - v f(u)
\neq 0$: sehingga determinan vektor $(u, f(u))$, $(v,
f(v))$ taknol, dan keduanya merentang $\R^2$.

\textbf{11.} Untuk $r, s \in \Q$: $f(ru + sv) = r f(u) + s f(v)$
(menurut pertanyaan 1 dua kali ditambah keaditifannya), sehingga grafiknya memuat
\[
\bigl(ru + sv,\; r f(u) + s f(v)\bigr)
= r\,(u, f(u)) + s\,(v, f(v)) .
\]
Diberikan $(x_0, y_0)$ dan $\varepsilon > 0$: sistem $2 \times 2$
yaitu $a(u, f(u)) + b(v, f(v)) = (x_0, y_0)$ mempunyai penyelesaian real
(yang tunggal) $(a, b)$ karena determinannya taknol. Lalu pilihlah
bilangan rasional $r_n \to a$, $s_n \to b$: maka $r_n(u, f(u)) +
s_n(v, f(v)) \to (x_0, y_0)$ secara koordinat demi koordinat, dan masing-masing
titiknya terletak pada grafiknya: jadi grafiknya \omterm{def:b1:topology:dense}{padat} di $\R^2$.

\textbf{12.} Misalkan $I$ sebuah \omterm{prop:b1:reals:intervals}{selang} yang taktermerosot, $x_0$ titik
tengahnya, dan $M$ sebarang: maka \omterm{def:b1:topology:dense}{kepadatannya} menyediakan sebuah titik grafik dalam jarak
$\min\bigl(\frac{\abs I}{2}, 1\bigr)$ dari $(x_0, M + 1)$, yakni
$x \in I$ dengan $f(x) > M$: jadi tak terbatas pada $I$, sehingga (menurut pertanyaan 8)
takkontinu di setiap titik dan tak monoton pada \omterm{prop:b1:reals:intervals}{selang} mana pun; dan
untuk sebarang sasaran $y_0$, titik grafik di dekat $(x_0, y_0)$ memberikan $f(x)$
yang sedekat-dekatnya dengan $y_0$ dengan $x \in I$: jadi $f(I)$ \omterm{def:b1:topology:dense}{padat} di
$\R$. Inilah pertanyaan 8 yang dibaca mundur: karena semua anak tangganya
setara, fungsi aditif yang taklinear pasti menggagalkan \emph{semuanya},
di mana-mana.

\textbf{13.} Keterdefinisiannya: karena $m + n\sqrt2 = m' + n'\sqrt2$ memaksa
$(n - n')\sqrt2 = m' - m \in \Z$, sehingga $n = n'$ (kalau tidak $\sqrt2 \in
\Q$) dan $m = m'$. Lalu keaditifannya pada $G$ jelas koordinat demi
koordinat. Ketakterbatasannya di dekat $0^+$: tetapkanlah $\varepsilon \in
\intoo{0}{1}$ dan $N \in \N$. Untuk setiap $n$ yang tetap dengan $\abs n
\leq N$, syarat $m + n\sqrt2 \in \intoo{0}{1}$ memakukan $m$
di dalam \omterm{prop:b1:reals:intervals}{selang} berpanjang $1$: jadi paling banyak satu bilangan bulat $m$ per
$n$, sehingga paling banyak $2N + 1$ unsur $G \cap \intoo{0}{1}$ yang mempunyai
$\abs{f} \leq N$. Padahal $G \cap \intoo{0}{\varepsilon}$
tak hingga (karena $G$ \omterm{def:b1:topology:dense}{padat}, \cref{exo:b1:reals:9}); sehingga ia
memuat suatu $g$ dengan $\abs{f(g)} > N$: jadi $f$ tak terbatas pada
setiap \omterm{def:b1:topology:open}{persekitaran} kanan $0$. Adapun memperluas $f$ menjadi fungsi
aditif pada $\R$ berarti memilih nilainya secara padu pada sebuah keluarga
bilangan real yang bebas linear atas $\Q$ dan merentang $\R$ atas $\Q$
--- yaitu basis Hamel; dan menghasilkan satu pun menuntut aksioma pemilihan,
serta tak ada rumus eksplisit yang dapat melakukannya: jadi secara konstruktif kita memiliki
monsternya hanya pada $G$.

\textbf{14.} Di sini $f(x) = f(\frac x2 + \frac x2) = f(\frac x2)^2 \geq
0$. Jika $f(x_0) = 0$, maka $f(x) = f(x - x_0)f(x_0) = 0$ untuk setiap
$x$: yang tersingkirkan. Jadi $f > 0$ dan $g = \ln \circ f$ bersifat \omterm{def:b1:continuity:continuous}{kontinu}
(\cref{prop:b1:continuity:algebra}) dengan $g(x + y) = g(x) +
g(y)$: sehingga menurut \cref{exo:b1:continuity:11}, $g(x) = cx$, jadi $f(x) =
\eu^{cx}$. Sebaliknya setiap $\eu^{cx}$ berlaku: jadi morfisme
\omterm{def:b1:continuity:continuous}{kontinu} $(\R, +) \to (\R^*, \times)$ tepat merupakan
eksponensialnya.

\textbf{15.} Fungsi $g(u) = f(\eu^u)$ \omterm{def:b1:continuity:continuous}{kontinu} dan $g(u + v) =
f(\eu^u \eu^v) = g(u) + g(v)$: sehingga $g(u) = cu$, dan setiap $x > 0$
tertulis $x = \eu^u$ dengan $u = \ln x$: jadi $f(x) = c\ln x$.

\textbf{16.} Fungsi $h = \ln \circ f$ \omterm{def:b1:continuity:continuous}{kontinu} pada
$\intoo{0}{+\infty}$ dengan $h(xy) = h(x) + h(y)$: sehingga menurut pertanyaan 15,
$h(x) = c\ln x$, jadi $f(x) = \eu^{c\ln x} = x^c$.

\textbf{17.} Dengan $x = y = 0$: $f(0) = 2f(0) + f(0)^2$, sehingga $f(0)(1 +
f(0)) = 0$. Jika $f(0) = -1$: maka menetapkan $y = 0$, $f(x) = f(x) +
f(0) + f(x)f(0) = -1$ untuk setiap $x$: yaitu konstanta $f \equiv -1$
(yang memang memenuhi persamaannya). Kalau tidak $f(0) = 0$; lalu $h
= 1 + f$ \omterm{def:b1:continuity:continuous}{kontinu}, $h(0) = 1$, dan
\[
h(x + y) = 1 + f(x) + f(y) + f(x)f(y) = h(x)\,h(y) :
\]
sehingga menurut pertanyaan 14, $h(x) = \eu^{cx}$, yakni $f(x) = \eu^{cx} - 1$
(dengan kasus $c = 0$ memberikan $f \equiv 0$). Jadi daftar lengkapnya: $f
\equiv -1$ dan $f(x) = \eu^{cx} - 1$, dengan $c \in \R$.

\textbf{18.} Dengan $x = y = 0$: $2f(0) = 4f(0)$, sehingga $f(0) = 0$. Dengan $x =
0$: $f(y) + f(-y) = 2f(y)$, sehingga $f$ genap. Dengan $y = x$: $f(2x) =
4f(x)$. Lalu induksi memakai $(x, y) \to (nx, x)$:
\[
f((n{+}1)x) = 2f(nx) + 2f(x) - f((n{-}1)x)
= (2n^2 + 2 - (n-1)^2) f(x) = (n+1)^2 f(x) .
\]
Lalu $f(x) = f\bigl(q\cdot\frac xq\bigr) = q^2 f\bigl(\frac
xq\bigr)$ memberikan $f\bigl(\frac pq x\bigr) = \frac{p^2}{q^2}f(x)$:
jadi $f(r) = r^2 f(1)$ pada $\Q$ (dengan kegenapannya menangani tandanya). Adapun kedua
fungsi \omterm{def:b1:continuity:continuous}{kontinu} $f$ dan $x \mapsto f(1)x^2$ bersesuaian pada
\omterm{def:b1:logic:sets}{himpunan} \omterm{def:b1:topology:dense}{padat} $\Q$, sehingga di mana-mana (pertanyaan 24): jadi $f(x) =
f(1)\,x^2$; dan setiap $cx^2$ memenuhi persamaannya.

\textbf{19.} Fungsi $g = f - f(0)$ \omterm{def:b1:continuity:continuous}{kontinu}, $g(0) = 0$, dan
memenuhi persamaan Jensen (karena konstantanya saling meniadakan). Dengan mengambil $y =
0$: $g\bigl(\frac x2\bigr) = \frac{g(x)}{2}$. Lalu untuk setiap $x,
y$:
\[
\frac{g(x + y)}{2} = g\Bigl(\frac{x+y}{2}\Bigr)
= \frac{g(x) + g(y)}{2} ,
\]
sehingga $g$ aditif dan \omterm{def:b1:continuity:continuous}{kontinu}: jadi $g(x) = cx$ (pertanyaan 2), dan
$f(x) = cx + d$ dengan $d = f(0)$. Adapun semua fungsi afin memenuhi
persamaan Jensen: sehingga daftarnya lengkap.

\textbf{20.} Lewat induksi pada $n$. Untuk $n = 0$: $\lambda \in \{0,
1\}$, yang sepele. Anggaplah ketaksamaannya berlaku untuk setiap bobot
$\frac{k}{2^n}$. Adapun bobot $\lambda = \frac{k}{2^{n+1}}$ dengan $k$ yang
genap tereduksi ke tingkat $n$; sedangkan untuk $k = 2j + 1$, $\lambda$ merupakan
titik tengah $\lambda_1 = \frac{j}{2^n}$ dan $\lambda_2 =
\frac{j+1}{2^n}$. Dengan $z_i = \lambda_i x + (1 - \lambda_i)y$:
$\lambda x + (1-\lambda)y = \frac{z_1 + z_2}{2}$, sehingga
\[
f\bigl(\lambda x + (1{-}\lambda)y\bigr)
\leq \frac{f(z_1) + f(z_2)}{2}
\leq \frac{(\lambda_1 + \lambda_2)f(x) + (2 - \lambda_1 -
\lambda_2)f(y)}{2}
= \lambda f(x) + (1 - \lambda)f(y) .
\]

\textbf{21.} Tetapkanlah $x, y$. Peta $\lambda \mapsto f(\lambda x +
(1 - \lambda)y)$ dan $\lambda \mapsto \lambda f(x) + (1 -
\lambda)f(y)$ bersifat \omterm{def:b1:continuity:continuous}{kontinu} pada $\intcc{0}{1}$ (menurut komposisi dan
aljabarnya, \cref{prop:b1:continuity:algebra}). Ketaksamaannya
berlaku pada bobot diadiknya, yang \omterm{def:b1:topology:dense}{padat} di $\intcc{0}{1}$
(\cref{exo:b1:reals:8}); lalu untuk $\lambda$ yang sebarang ambillah bilangan diadik
$\lambda_n \to \lambda$ lalu beralihlah ke limitnya
(\cref{thm:b1:continuity:seqchar} dan \cref{thm:b1:seq:order}):
sehingga ketaksamaan kecembungannya berlaku untuk setiap $\lambda \in
\intcc{0}{1}$ --- jadi kecembungan titik tengah ditambah \omterm{def:b1:continuity:continuous}{kekontinuan} sama dengan
kecembungan (yaitu gagasan pada \cref{ch:b1:derivative}).

\textbf{22.} Fungsi aditif $f$ yang taklinear memenuhi
$f\bigl(\frac{x+y}{2}\bigr) = \frac{f(x) + f(y)}{2}$ dengan persis
(menurut pertanyaan 1 dengan $r = \frac12$, lalu keaditifannya): jadi ia
cembung titik tengah, bahkan afin titik tengah. Seandainya ia memenuhi ketaksamaan
kecembungan yang penuh pada suatu \omterm{prop:b1:reals:intervals}{selang} $\intcc{x}{y}$, maka untuk
$\lambda \in \intcc{0}{1}$: $f(\lambda x + (1-\lambda)y) \leq
\max(f(x), f(y))$ --- yakni terbatas di atas pada \omterm{prop:b1:reals:intervals}{selang} yang taktermerosot,
sehingga linear menurut pertanyaan 9: yang bertentangan. Jadi \omterm{def:b1:continuity:continuous}{kekontinuan} pada
pertanyaan 21 bukanlah kemewahan: karena tanpanya, kecembungan titik tengah
hanya mengendalikan kerangka diadik yang terbilang, sedangkan kontinum
di antaranya berlarian liar.

\textbf{23.} Fungsi $p(x) = \frac{x^2}{2}$ memenuhi $p(x+y) = p(x) +
p(y) + xy$. Jika $f$ sebarang penyelesaian yang \omterm{def:b1:continuity:continuous}{kontinu}, maka $g = f - p$
\omterm{def:b1:continuity:continuous}{kontinu} dan aditif, sehingga $g(x) = cx$:
\[
f(x) = \frac{x^2}{2} + cx , \qquad c \in \R ,
\]
dan masing-masingnya merupakan penyelesaian: jadi daftarnya lengkap.

\textbf{24.} Asas umumnya: jika $u, v$ \omterm{def:b1:continuity:continuous}{kontinu} dan
bersesuaian pada $D \subseteq \R$ yang \omterm{def:b1:topology:dense}{padat}, maka untuk $x \in \R$ pilihlah
$d_n \in D$ dengan $d_n \to x$
(\cref{prop:b1:topology:closureprops}); sehingga $u(x) = \lim u(d_n) =
\lim v(d_n) = v(x)$. Sekarang misalkan $f$ \omterm{def:b1:continuity:continuous}{kontinu} dan aditif pada
\omterm{def:b1:structures:subgroup}{subgrup} \omterm{def:b1:topology:dense}{padat} $G$. Tetapkanlah $x, y \in \R$ lalu ambillah $g_n \to x$,
$g'_n \to y$ dengan $g_n, g'_n \in G$; maka $g_n + g'_n \to x + y$
dan, menurut \omterm{def:b1:continuity:continuous}{kekontinuan} barisannya di $x + y$, $x$ dan $y$:
\[
f(x + y) = \lim f(g_n + g'_n) = \lim\bigl(f(g_n) +
f(g'_n)\bigr) = f(x) + f(y) :
\]
jadi $f$ aditif pada seluruh $\R$ (sehingga linear, menurut pertanyaan 2).

\textbf{25.} (i) Untuk $f$ yang aditif: linear $\iff$ \omterm{def:b1:continuity:continuous}{kontinu}
$\iff$ \omterm{def:b1:continuity:continuous}{kontinu} di satu titik $\iff$ monoton pada suatu \omterm{prop:b1:reals:intervals}{selang}
$\iff$ terbatas (bahkan secara sepihak) pada suatu \omterm{prop:b1:reals:intervals}{selang}. (ii) \omterm{def:b1:topology:dense}{Kepadatan}
grafiknya di bidangnya memperlihatkan bahwa kegagalannya bukan cacat
lokal melainkan ledakan yang menyeluruh: karena di atas setiap subselangnya
nilainya melumur ke seluruh $\R$, sehingga setiap sifat keteraturannya
gagal di mana-mana serentak. (iii) Tangkapannya: $cx$, $\eu^{cx}$,
$c\ln x$, $x^c$, $cx^2$, dan $cx + d$ --- yaitu enam pencirian,
dengan satu metode: angkutlah persamaannya ke persamaan Cauchy, lalu buktikan
kerangka $\Q$-nya lewat induksi, lalu naikkan pangkatnya ke $\R$ lewat \omterm{def:b1:topology:dense}{kepadatan} ditambah
\omterm{def:b1:continuity:continuous}{kekontinuan}. (iv) Adapun \omterm{def:b1:topology:dense}{kepadatan} $\Q$ (atau bilangan diadik) memikul
kenaikan pangkat pada \cref{exo:b1:continuity:11} dan pada pertanyaan 21; sedangkan ia
tak memikul apa pun pada pertanyaan 13, karena tanpa \omterm{def:b1:continuity:continuous}{kekontinuan}
nilainya tak merambat dari \omterm{def:b1:logic:sets}{himpunan} yang \omterm{def:b1:topology:dense}{padat} ke \omterm{def:b1:topology:closure}{penutupnya} ---
jadi \omterm{def:b1:topology:dense}{kepadatan} memindahkan informasi hanya sepanjang \omterm{def:b1:continuity:continuous}{kekontinuan}.

\end{solution}
